% DRAFT 4 — Related Work. Former 2.1 (130w) + 2.2 (75w) merged per Law 6.
% Reframed: the gap is now BOTH geographic and methodological.

\section{Related Work}
\label{sec:related}

\paragraph{Financial and regulatory benchmarks.}
The financial NLP community has built substantial evaluation infrastructure, anchored almost
entirely in Western text. FinQA \citep{chen2021finqa} and ConvFinQA \citep{zheng2022convfinqa} test
numerical reasoning over SEC filings; FiNER-139 \citep{loukas2022finer} targets entity recognition
in the same filings; FLUE \citep{shah2022flue} packages several such tasks into a multi-task suite.
PIXIU/FLARE \citep{xie2023pixiu} and FinBen \citep{xie2024finben} push toward breadth, the latter
spanning 42 datasets across 24 tasks. FinanceBench \citep{islam2023financebench} is closest to our
own construction---expert-verified gold answers over financial documents---but covers publicly
listed US companies only. Legal NLP contributes adjacent structure without the financial content:
CUAD \citep{hendrycks2021cuad} for contract clauses, LexGLUE \citep{chalkidis2022lexglue} for
European legal text, ILDC \citep{malik2021ildc} for Indian court judgments. Across this literature
the binding constraint is geography rather than task diversity, and no benchmark we are aware of
defines amendment-chain supersession as a task type.

\paragraph{How benchmarks are scored.}
The methodological gap matters more to this paper. Extractive QA evaluation inherits its scoring
conventions from SQuAD \citep{rajpurkar2016squad} and DROP \citep{dua2019drop}: exact match and
token-overlap F1 against a reference string. Those conventions were designed for short spans copied
from a passage, where the space of correct surface forms is small. Applied to open-ended
quantitative answers---where \texttt{Rs.\ 19,000 crore}, \texttt{19,000 crore}, and
\texttt{nineteen thousand crore} are the same answer---they measure formatting as much as
correctness. LLM-as-judge evaluation has emerged as the usual alternative, and its biases are by now
well documented, including a tendency to over-credit fluent or elaborated reasoning. Our position is
not that one regime is correct. It is that the two regimes answer different questions, that the
difference is large enough to reorder a leaderboard, and that reporting either one alone leaves the
reader unable to tell capability from format compliance. Holistic frameworks such as HELM
\citep{liang2022helm} argue for multi-metric reporting on related grounds; we supply a concrete case
where the metric choice, not the metric count, is what moves the ranking.

\paragraph{What a benchmark reports about itself.}
Benchmarks are commonly described by nominal item count. Whether those items discriminate between
the systems being ranked is rarely stated, though the psychometrics literature has long treated item
difficulty and discrimination as basic descriptive properties of a test. Section~\ref{sec:effsize}
reports them for IndiaFinBench.
